\documentclass[11pt]{article}
\usepackage[margin=2.5cm]{geometry}
\usepackage{amsmath,amssymb,amsthm}
\usepackage{hyperref}

\newtheorem{lemma}{Lemma}
\newtheorem{proposition}{Proposition}
\theoremstyle{remark}
\newtheorem*{remark}{Remark}
\newtheorem*{example}{Example}

\newcommand{\E}{\mathbb{E}}
\newcommand{\Cov}{\operatorname{Cov}}
\newcommand{\tr}{\operatorname{tr}}
\newcommand{\ts}{\theta^*}
\newcommand{\hth}{\hat\theta_m}

\title{The plug-in learner coefficient $\tfrac12\tr(FH^{-1})$:\\ sandwich covariance, excess loss and the information equality}
\author{Notes on Section 3.7 of the Stage 1 write-up}
\date{21 October 2026}

\begin{document}
\maketitle

\noindent The write-up predicts the regret slope of a plug-in (point-estimate) code by $\kappa=\tfrac12\tr(FH^{-1})$, with $H$ the Hessian of the expected loss and $F$ the second moment of the per-example gradients at the optimum. These notes derive that formula in three steps (Sections~\ref{sec:sandwich}--\ref{sec:sum}), prove that $F=H$ for a correctly specified log-likelihood (Section~\ref{sec:infoeq}), and list where each step can fail (Section~\ref{sec:fail}).

\section{Setup}

Data $z_1,z_2,\dots$ are i.i.d., $z=(x,y)$. The per-example loss is the negative log-likelihood of a conditional model,
\[
\ell(z,\theta)=-\log p_\theta(y\mid x),\qquad \theta\in\mathbb R^d .
\]
The population risk and its minimiser are
\[
L(\theta)=\E\,\ell(z,\theta),\qquad \ts=\arg\min_\theta L(\theta),\qquad \nabla L(\ts)=0,
\]
and the empirical risk on $m$ examples and its minimiser (the plug-in fit) are
\[
L_m(\theta)=\frac1m\sum_{i=1}^m\ell(z_i,\theta),\qquad \hth=\arg\min_\theta L_m(\theta),\qquad \nabla L_m(\hth)=0 .
\]
The two matrices, both evaluated at $\ts$:
\[
H=\nabla^2L(\ts)=\E\,\nabla^2\ell(z,\ts),\qquad
F=\E\bigl[\nabla\ell(z,\ts)\,\nabla\ell(z,\ts)^{\top}\bigr].
\]
$H$ is the \emph{curvature} of the expected loss: how fast it rises as $\theta$ leaves $\ts$. $F$ is the \emph{spread} of the individual examples' gradients: because $\E\,\nabla\ell(z,\ts)=\nabla L(\ts)=0$, it is their covariance, i.e.\ how much single examples disagree about which way to move. Both are computed from the same data distribution; they are different functionals of it, not the same quantity on different subsets.

\section{The sandwich covariance}
\label{sec:sandwich}

\begin{lemma}[Covariance under a linear map]\label{lem:lin}
If $g$ is a random vector with mean $0$ and covariance $\Sigma$, and $A$ is a fixed matrix, then $\Cov(Ag)=A\Sigma A^\top$.
\end{lemma}
\begin{proof}
$\Cov(Ag)=\E[Ag\,(Ag)^\top]=A\,\E[gg^\top]\,A^\top=A\Sigma A^\top$. The matrix appears twice because a covariance is quadratic in the vector.
\end{proof}

\begin{proposition}[Sandwich]\label{prop:sandwich}
Under the regularity conditions of Section~\ref{sec:fail} (in particular $H$ positive definite),
\[
\sqrt m\,(\hth-\ts)\ \xrightarrow{d}\ \mathcal N\bigl(0,\;H^{-1}FH^{-1}\bigr).
\]
\end{proposition}
\begin{proof}[Derivation]
Expand the first-order condition $\nabla L_m(\hth)=0$ around $\ts$ (mean-value form, $\bar\theta$ between $\hth$ and $\ts$):
\[
0=\nabla L_m(\ts)+\nabla^2L_m(\bar\theta)\,(\hth-\ts)
\quad\Longrightarrow\quad
\hth-\ts=-\bigl[\nabla^2L_m(\bar\theta)\bigr]^{-1}\,\nabla L_m(\ts).
\]
Two limit theorems handle the two factors.
\begin{itemize}
\item \emph{Law of large numbers.} $\hth\to\ts$ (consistency), hence $\bar\theta\to\ts$, and the empirical Hessian converges to its mean: $\nabla^2L_m(\bar\theta)\to H$.
\item \emph{Central limit theorem.} $\nabla L_m(\ts)=\frac1m\sum_i\nabla\ell(z_i,\ts)$ is an average of i.i.d.\ vectors with mean $0$ and covariance $F$, so $\sqrt m\,\nabla L_m(\ts)\xrightarrow{d}\mathcal N(0,F)$.
\end{itemize}
By Slutsky's lemma, $\sqrt m(\hth-\ts)\approx -H^{-1}\cdot\sqrt m\,\nabla L_m(\ts)$, and Lemma~\ref{lem:lin} with $A=-H^{-1}$ (symmetric) gives the covariance $H^{-1}FH^{-1}$.
\end{proof}

\begin{remark}[One dimension]
With scalars $h=L''(\ts)$ and $f=\E[\ell'(z,\ts)^2]$, the derivation reads $\hth-\ts\approx-g/h$ with $g$ the average gradient, $\operatorname{Var}(g)=f/m$, so
\[
\operatorname{Var}(\hth)\approx\frac{f}{m\,h^2}.
\]
Noisy per-example gradients (large $f$) make the fit wander; a sharply curved loss (large $h$) pins it down. The name ``sandwich'' only describes the shape $H^{-1}\,F\,H^{-1}$: $F$ between two slices of $H^{-1}$.
\end{remark}

\section{The excess loss of the fit}
\label{sec:excess}

The quantity a prequential code pays at step $m$ is the expected loss of the fit $\hth$ on a fresh example, in excess of the best predictor's:
\[
\E\bigl[L(\hth)\bigr]-L(\ts).
\]
Expand $L$ around $\ts$; the linear term vanishes because $\nabla L(\ts)=0$:
\[
L(\hth)-L(\ts)=\tfrac12(\hth-\ts)^\top H(\hth-\ts)+o\bigl(\|\hth-\ts\|^2\bigr).
\]
For any random vector $u$ with mean $0$, $\E[u^\top Hu]=\E\tr(Huu^\top)=\tr(H\Cov u)$. With $\Cov(\hth-\ts)\approx H^{-1}FH^{-1}/m$ from Proposition~\ref{prop:sandwich}:
\[
\E\bigl[L(\hth)\bigr]-L(\ts)\approx\frac{1}{2m}\tr\bigl(H\,H^{-1}FH^{-1}\bigr)=\frac{\tr(FH^{-1})}{2m}=\frac{\kappa}{m},
\qquad
\boxed{\ \kappa=\tfrac12\tr(FH^{-1})\ }.
\]
In one dimension: $\tfrac12 h\cdot f/(mh^2)=f/(2mh)$. One factor of $h$ cancels; what remains is the ratio of gradient noise to curvature.

\begin{remark}
Passing from the distributional limit to the expectation needs uniform integrability of $m\|\hth-\ts\|^2$. This is the standard heuristic and holds in regular models; it is where heavy tails or rare divergent fits would break the formula.
\end{remark}

\section{From excess loss to the regret slope}
\label{sec:sum}

A per-sample prequential code codes target $m+1$ with the fit on the first $m$ targets. Summing the expected excess,
\[
\E\,\mathrm{Reg}(n)\approx\sum_{m=1}^{n}\frac{\kappa}{m}=\kappa\ln n+O(1),
\]
so $\kappa$ is the slope of the regret against $\ln n$. (Block codes with geometric ratio $r$ multiply this slope by $(r-1)/\ln r$; the write-up divides it out as the ``batching correction''.)

\section{The information equality: \texorpdfstring{$F=H$}{F=H} under correct specification}
\label{sec:infoeq}

Let $s(\theta)=\nabla_\theta\log p_\theta(y\mid x)$ be the score. Then
\[
\nabla\ell=-s,\qquad
\nabla^2\ell=-\nabla\frac{\nabla p_\theta}{p_\theta}=-\frac{\nabla^2p_\theta}{p_\theta}+\frac{\nabla p_\theta\,\nabla p_\theta^\top}{p_\theta^2}
=-\frac{\nabla^2p_\theta}{p_\theta}+s\,s^\top .
\]
Taking expectations at $\ts$:
\begin{equation}
H=F-\E\!\left[\frac{\nabla^2p_{\ts}(y\mid x)}{p_{\ts}(y\mid x)}\right].
\label{eq:HF}
\end{equation}

\begin{proposition}[Information equality]
If the model is correctly specified at $\ts$, i.e.\ the true conditional law of $y$ given $x$ is $p_{\ts}(\cdot\mid x)$ for (almost) every $x$, then $H=F$, both equal to the Fisher information $I(\ts)=\E[s s^\top]$.
\end{proposition}
\begin{proof}
Condition on $x$ and take the expectation over $y$ under the true law, which is $p_{\ts}$ itself. The density cancels against the denominator in~\eqref{eq:HF}:
\[
\E_{y\sim p_{\ts}(\cdot\mid x)}\!\left[\frac{\nabla^2p_{\ts}(y\mid x)}{p_{\ts}(y\mid x)}\right]
=\sum_y\nabla^2p_\theta(y\mid x)\Big|_{\ts}
=\nabla^2\Bigl(\sum_yp_\theta(y\mid x)\Bigr)\Big|_{\ts}
=\nabla^2 1=0 .
\]
Averaging over $x$ (under any distribution) gives $H=F$.
\end{proof}

\noindent Consequently $\tr(FH^{-1})=\tr(I_d)=d$ and $\kappa=d/2$: the classical cost of $\tfrac12\ln n$ per parameter, which in regular models coincides with the Bayes learning coefficient $\lambda=d/2$.

\paragraph{Assumptions of the information equality.}
\begin{enumerate}
\item \emph{Correct specification at $\ts$.} The conditional law of $y$ given $x$ is exactly $p_{\ts}(\cdot\mid x)$. The distribution of $x$ is arbitrary; only the $y$-expectation must be under the model's own law. This is the cancellation step.
\item \emph{The loss is the log-likelihood.} No temperature, reweighting or penalty: for $\ell=-\tau\log p_\theta$, or with a regulariser added, $F\neq H$ even when the model is correct.
\item \emph{Differentiation and summation commute,} and the support of $p_\theta$ does not depend on $\theta$. Automatic for finitely many classes; for continuous $y$ it needs dominated convergence.
\item \emph{Smoothness and positivity.} $p_\theta$ is twice differentiable near $\ts$ and positive wherever the true law has mass (so that $\log p$ and $s$ are defined).
\item \emph{Only at $\ts$.} Away from the optimum, or for the empirical versions at finite $m$, $F$ and $H$ differ; they agree only in expectation at $\ts$, and the empirical estimates only asymptotically.
\end{enumerate}

\begin{example}[Misspecified noise scale]
Model $y\sim\mathcal N(\theta,\sigma^2)$, truth $y\sim\mathcal N(\mu,\tau^2)$. Then $\ell=(y-\theta)^2/(2\sigma^2)+\text{const}$, $\ts=\mu$, and
\[
H=\frac1{\sigma^2},\qquad F=\E\frac{(y-\mu)^2}{\sigma^4}=\frac{\tau^2}{\sigma^4},\qquad \kappa=\frac{\tau^2}{2\sigma^2}.
\]
The mean is estimated just as well either way, but an overconfident model ($\sigma<\tau$) pays more per $\ln n$ than $d/2=\tfrac12$, and an underconfident one less. Correct specification $\sigma=\tau$ restores $F=H$ and $\kappa=\tfrac12$. Coding randomized-response labels with the clean softmax is the classification analogue of $\sigma<\tau$.
\end{example}

\begin{remark}[Empirical versus model Fisher]
The write-up's $F$ is the \emph{empirical} Fisher: gradients evaluated at the observed labels. That is the $F$ of the sandwich, the right quantity for $\kappa$. The \emph{model} Fisher $\E_x\E_{y\sim p_{\ts}}[ss^\top]$ takes the $y$-expectation under the model; it equals $F$ only under correct specification.
\end{remark}

\section{Where the steps fail}
\label{sec:fail}

\paragraph{Assumptions of the sandwich and excess-loss steps} (needed even when $F\neq H$): i.i.d.\ data; a unique optimum $\ts$ in the interior of the parameter set, to which $\hth$ converges; $\ell$ twice continuously differentiable with dominated derivatives near $\ts$; $F$ finite; and $H$ positive definite.

\paragraph{Singular models.} If $H$ is singular, $H^{-1}$ does not exist: along the flat directions the fit does not concentrate at rate $m^{-1/2}$, the quadratic expansion is not the leading term, and Proposition~\ref{prop:sandwich} fails. This is the regime of Watanabe's theory, in which the Bayes learning coefficient satisfies $\lambda<d/2$ and a plug-in's coefficient need not equal $\lambda$. Exact gauge directions (such as $W\to WA$ for the subspace map) are harmless if they are quotiented out; the problem is directions that are flat to leading order but not exactly.

\paragraph{The write-up's networks.} On GPT-2 the formula is applied with a pseudo-inverse on the non-flat directions. There $F\neq H$ for two reasons: the target law is only approximately realizable (it is reached only at infinite logits, so assumption~1 holds only in the limit), and the readout is saturated, so the per-example gradients are tiny in most directions. Few directions carry most of $F$, and $\tfrac12\tr(FH^{-1})$ comes out two orders of magnitude below the family's nominal $\lambda=d_V(d-d_V)/2$: the code pays only for the directions the data actually pull on.

\end{document}
